% 学習結果説明（記号表。すべての説明書で使う記号・モデル名・列名の一覧）
% results フォルダで「lualatex 学習結果説明_記号表.tex」を2回実行すると PDF を作り直せる（TeX Live の LuaLaTeX を使用）。
% 04_補助スクリプト/build_report_symbols.py が、report_symbols_template.tex から作る。新しい記号・モデル名・列名を使ったら、ここにも足す（CLAUDE.md §8）。
\documentclass[paper=a4paper,fontsize=10.5pt,line_length=44zw,number_of_lines=38]{jlreq}
\usepackage[haranoaji]{luatexja-preset}
\usepackage{amsmath}
\usepackage{booktabs}
\usepackage{array}
\usepackage{longtable}
\usepackage{enumitem}
\usepackage{xcolor}
\usepackage[hidelinks,bookmarksnumbered]{hyperref}

\newcommand{\file}[1]{\texttt{\detokenize{#1}}}
\newcommand{\gd}{\dot{\gamma}}
\newcommand{\Tff}{T_{\mathrm{ff}}}
\newcommand{\Tset}{T_{\mathrm{set}}}
\newcommand{\Eh}{E_{\mathrm{h}}}
\newcommand{\Ev}{E_{\mathrm{v}}}
\newcommand{\PI}{P3\_1}
\newcommand{\PII}{P3\_2}
\newenvironment{points}{\begin{itemize}[label=$\triangleright$,leftmargin=1.5\zw,itemsep=0.3ex,topsep=0.5ex,parsep=0pt]}{\end{itemize}}
\newcommand{\pt}[1]{\item \textbf{#1}：}
\renewcommand{\arraystretch}{1.25}
\setlength{\LTpre}{0.6ex}
\setlength{\LTpost}{1.2ex}

% 記号表の表：記号／呼び名と意味／単位・値／出てくる説明書（2・3a・3b・3c・3d・5a。全：すべて）
\newenvironment{symtab}{%
  \begin{longtable}{@{}>{\raggedright\arraybackslash}p{6.5\zw}>{\raggedright\arraybackslash}p{22\zw}>{\raggedright\arraybackslash}p{6\zw}>{\raggedright\arraybackslash}p{4.5\zw}@{}}
  \toprule 記号 & 呼び名と意味 & 単位・値 & 出る所 \\ \midrule \endfirsthead
  \toprule 記号 & 呼び名と意味 & 単位・値 & 出る所 \\ \midrule \endhead
  \bottomrule \endlastfoot}{\end{longtable}}

\title{樹脂粘度モデル開発　記号表\\[0.6ex]{\large すべての学習結果説明（手順2・3a〜3e・5a〜5b・材料別 GPR）で使う記号・モデル名・列名の一覧}}
\author{}
\date{@@DATE@@}

\begin{document}
\maketitle
{\small\tableofcontents}

% =====================================================================
\section{この文書について}

各手順の学習結果説明（\file{学習結果説明.pdf}・\file{学習結果説明_手順3a.pdf}・\file{学習結果説明_手順3d.pdf} など）には、それぞれに用語集があります。この文書は、\textbf{それらで説明済みの記号も含めて、すべての記号・モデル名・データの呼び名・ノートと CSV の列名を、1か所にまとめた一覧}です。説明書を読んでいて記号が分からなくなったら、ここで引いてください。

\begin{points}
\pt{表の形}「記号」「呼び名と意味」「単位・値」「出る所」の4列です。「出る所」は、その記号が出てくる説明書で、2＝手順2（\file{学習結果説明.pdf}）、3a・3b・3c・3d・3e・5a＝各手順の説明書、材＝材料別 GPR（05 フォルダ。別のプロジェクト）、厚＝その厚みごとの比較（\file{学習結果説明_厚みごとの温度シフト.pdf}）、全＝すべてです。
\pt{先に読む所}同じ文字が、手順によって別の意味で使われています（H、A、B、C、F、L、R、S、T、P、E など）。2章の「紛らわしい記号」を、先に読んでください。
\pt{決まり}$\ln$ は自然対数、$\log_{10}$ は常用対数です。目的変数は、すべて $\ln\eta$（粘度の自然対数）です。
\pt{更新のしかた}新しい記号・モデル名・列名を使ったときは、この文書にも足します（CLAUDE.md §8）。\file{results} フォルダで \texttt{lualatex 学習結果説明\_記号表.tex} を2回実行すると、PDF を作り直せます。
\end{points}

% =====================================================================
\section{紛らわしい記号（先に読む）}

同じ文字が、文脈で別の意味になるものを並べます。「出る所」の違いを見てください。

\begin{symtab}
$H$ & キャビティ厚み（斜体）。\textbf{構造 H}（5a の当初の温度シフト構造。「構造 H」と書く）とは別物。 & 0.5〜2.0 mm & 全／5a \\
A & \textbf{構造 A}（5a。縦方向だけシフト）。3a の \textbf{A0・A1・A2・A1$'$}（圧力の有無を比べる GPR）の先頭の A は、構造とは無関係。 & ― & 5a／3a \\
B & \textbf{構造 B}（5a。横と縦を別々にシフト）。3b・3c の \textbf{B0・B0T など}（ショットごとの $\gd$ を使う GPR）の B は、構造とは無関係。 & ― & 5a／3b・3c \\
C & \textbf{構造 C}（5a。構造なしの通常の GPR）。カーネルの \textbf{ConstantKernel}（振幅$^2$）の C とも別。 & ― & 5a／全 \\
F & \textbf{構造 F}（5a。$\Tff$ でシフトする参考の構造）。モデル名の末尾の F（$\Tff$ あり：B0F、R3F）とも別。フーリエ数は $\mathrm{Fo}$（2 文字）。 & ― & 5a／3b \\
L & センサ間距離 $L$（17.5 mm）。モデル \textbf{L0・L1・L2・L0T}（線形回帰）の L は、Linear（線形）の頭文字。 & mm／― & 全／3a・3b \\
R & 気体定数 $R$（8.314 J/mol/K）。モデル \textbf{R0・R2・R3}（条件ごとの値を使う GPR。Reference）の R。決定係数 $R^2$。収束の指標 R-hat。厚みごとの比較では、温度による粘度の比 $R$。 & ― & 5a／3b・3c／全／厚 \\
S & データの分け方 \textbf{S\_V・S\_T・S\_R1}、手順2の \textbf{S1・S2}。小文字の $s$ は、流動曲線の傾き。 & ― & 全／5a \\
T & 手順2：旧データの流動先端温度 $T$。新データでは $\Tff$。5a の式の $T$：絶対温度（K）で、設定温度 $\Tset$（または $\Tff$）を K にした値。 & ℃／K & 2／5a \\
P & 手順2：旧データの樹脂圧力 $P$（Pa）。新データ：\PI・\PII（MPa）。モデル名 R3P の P は、\PII\ を使うことを表す。 & Pa／MPa & 2／3a・3b \\
E & 活性化エネルギー $E$（構造 H・F）。$\Eh$・$\Ev$（構造 B など）。$E$ は、文献の活性化エネルギーとは別物の「見かけのシフトエネルギー」。 & kJ/mol & 5a \\
M & 手順2の \textbf{M1〜M4}（モデル）。MIX は、粉砕材混合 PP の略号。 & ― & 2／全 \\
N & 構造 N（温度なし、$E=0$）。 & ― & 5a \\
R3 & CLAUDE.md §6 の「R1」（比較用の通常の GPR）と同じ入力。手順3b の R3、手順5a の C（入力の $\gd$ は条件平均）、手順3d の R3 も同じ形。 & ― & 3b〜5a \\
$z$ & 手順3d の規則 f の $z$（残差を予測標準偏差で割ったもの。しきい値 3.5）。外れ値なしのファイルを作ったときの「修正 $z$ スコア」（Tff に使う。しきい値 3.5）とは別。 & ― & 3d／3a \\
5data 版・Alldata 版 & 手順3d 以降の「データの版」。5data 版：5data で学習し、3data で外挿を検証する。Alldata 版：全条件（PP は 1 条件 10 ショット）で学習する。ファイル名の 5data・Alldata（ファイルそのもの）とは、別の言葉（版は、学習と検証の組み方）。 & ― & 3d \\
S\_V と S\_Vall & S\_V は 5data で学習して 3data を検証（5data 版）。S\_Vall は、Alldata のうち 5data の条件だけ（全ショット）で学習して 3data を検証（Alldata 版。PP だけ意味がある）。 & ― & 3d \\
\file{shot} と \file{pos} & \file{shot}：元ファイル名の先頭 3 桁（設定温度・厚みごとの通し番号）。\file{pos}：その条件の中でのショットの順番（最初 = 1）。ショットの順番によるずれは \file{pos} で調べる。 & ― & 3d \\
CV & 条件単位の交差検証（K 分割。\file{CV_K}、既定 10）。手順3d の Alldata 版の評価。LOCO（条件を 1 つずつ抜く）とは別（分割数が違う）。 & ― & 3d／3a・3b \\
\end{symtab}

% =====================================================================
\section{測定と物理量}

\begin{symtab}
$\eta$ & 見かけの粘度。せん断発熱・冷却・固化の影響を含む「見かけ」の値。目的変数は $\ln\eta$。 & Pa$\cdot$s & 全 \\
$\gd$ & 見かけのせん断速度。$\gd = 6L/(H\,\Delta t)$。ショットごとの値。入力には $\ln\gd$ を使う。 & 1/s & 全 \\
$\gd$（条件平均）& 同じ条件のショットで、$\ln\gd$ を平均した値（cm：condition mean）。成形条件から見たばらつきを出すモデル（R 系）の入力。 & 1/s & 3b・3c・5a \\
$\tau$、$\tau_w$ & 見かけの壁面せん断応力。$\tau = \eta\gd = \Delta P\,H/(2L)$。 & Pa & 2・3a \\
$\Delta P$ & 圧力損失。新データでは $\PI - 0.30$ MPa、旧データでは $P - P_2$。 & Pa & 全 \\
$L$ & センサ間距離。 & 17.5 mm & 全 \\
$H$ & キャビティ厚み。0.5・1.0・1.5・2.0 mm の4水準。入力には $\ln H$。 & mm & 全 \\
$\Delta t$ & 2つの圧力センサへの樹脂到達の時間差（出力の立ち上がり）。5 µs の倍数。 & ms（7〜1,170） & 3a・5a \\
$v$ & 射出速度（成形条件として決める設定値）。射出率 [cm$^3$/s] はこの 0.531 倍。入力には $\ln v$（モデル R2）。 & mm/s & 3a〜5a \\
$\Tset$ & 加熱シリンダ設定温度（成形条件として決める設定値）。モデル上の「溶融温度」。入力には ℃ のまま（標準化して）使う。 & 190・210・230 ℃ & 3a〜5a \\
$\Tff$ & 流動先端温度（流動中樹脂温度）。センサ間を流れる樹脂の先端温度の実測値。流れの結果で、事前には分からない。入力には $\ln\Tff$（K）。 & ℃ & 3a〜5a \\
$T$（手順2）& 旧データの流動先端温度。新データの $\Tff$ に当たる。 & ℃ & 2 \\
$P$（手順2）& 旧データの流動中樹脂圧力（1点目のセンサ）。 & Pa & 2 \\
$P_2$ & 旧データで、2点目のセンサの圧力にあたるオフセット（約 0.400 MPa）。 & MPa & 2 \\
\PI & 流動中樹脂圧力。センサ4に樹脂が届いた時刻のセンサ3の圧力。粘度の算出式 $\eta\gd = H(\PI - 0.30)/(2L)$ に入っているので、入力に使うと答えが計算できてしまう（リーク）。使わない。 & MPa（Pa に直して $\ln$） & 3a〜 \\
\PII & 流動中樹脂圧力。センサ4が最大値になった時刻のセンサ3の圧力。算出式には入っていないが、流れの結果で、事前には分からない。 & MPa（Pa に直して $\ln$） & 3a〜 \\
$c$（0.30 MPa）& 算出式のオフセット。旧データは 0.40 MPa。 & MPa & 3a \\
$s$ & 見かけの流動曲線の傾き $d\ln\eta/d\ln\gd$。$-1$ ならせん断応力が一定。新データの中央値は約 $-1.6$（旧データは $-0.88$）。 & ― & 5a \\
$\mathrm{Fo}$ & フーリエ数。冷却の度合い。$\mathrm{Fo} = \alpha\Delta t/H^2$。入力は $\ln\mathrm{Fo} = -(\ln\gd + 3\ln H)$ ＋ 定数。 & ― & 5a \\
$\alpha$ & 熱拡散率。定数なので GP には影響しない。 & ― & 5a \\
$\ln$、$\log_{10}$ & 自然対数（底 $e\approx2.718$）、常用対数（底 10）。$\ln\eta$ の差 0.1 は、粘度の約 1.1 倍の違い。 & ― & 全 \\
条件 & 設定温度・厚み・射出速度の組み合わせ（計 129 条件）。ノートでは \file{cond} 列（例：\file{190_0.5_20}）。 & ― & 全 \\
ショット & 同じ条件での1回の成形。元ファイル名の先頭 3 桁（\file{shot} 列）。この番号は、設定温度・厚みごとの通し番号で、条件の中での順番は \file{pos} 列。 & ― & 全 \\
ショットの順番によるずれ & 条件の中の $\ln\eta$（など）の、条件平均からの偏差を、順番 \file{pos} に回帰した傾き。「ずれ（最初→最後）」はその傾きに、順番の幅（PP は 9、粉砕材混合 PP は 4）を掛けたもの。許容値 \file{DRIFT_TOL} = 0.02。 & ln η & 3d \\
残差・予測標準偏差 & 規則 f：条件を 1 つずつ抜いた GPR の予測と、実測（条件平均）の差が「残差」（$\ln\eta$）。「予測標準偏差」は、その予測の不確かさ。$z$ = 残差 / 予測標準偏差。 & ln η & 3d \\
規則 f & 外れ値の検知の規則。$|z|>3.5$ かつ |残差| > 0.1（粘度で約 10\%）の条件を拾う。拾った条件は、問題データ一覧に「除外」で記録して除いてよい（ユーザー、2026-10-14）。規則 a〜e は、見直しの候補を拾う規則で、除くかはユーザーが決める。 & ― & 3d \\
標準誤差 & 推定した値（傾きなど）の、不確かさの目安。 & ― & 3d \\
純粋誤差 & 同じ条件のショット間のばらつき（$\ln\eta$ の標準偏差）。外れ値なしのデータは、20\% 以上離れた値を除いているので、小さめに出る。成形条件だけを入力にしたモデルは、これより小さい誤差にはなれない（ノイズの床）。リークがあると、これより小さくなる（手順3e の目印）。 & ― & 3a・3b・3e \\
理想的なスリット流れ & 等温で、固化層も壁面すべりもない流れ。見かけの流動曲線が、厚みによらず 1 本に重なる（Mooney の方法の前提）。重ならなければ、冷却・固化層などの効果がある。 & ― & 3e \\
固化層 & 金型の壁で冷えて固まった樹脂の層（片側の厚さ $\delta$）。流れられる隙間が $H - 2\delta$ に狭まる。 & mm & 3e \\
厚みによる粘度の比 & 同じ $\gd$・同じ設定温度での $\eta(H)/\eta(2.0\,\mathrm{mm})$（\file{thickness_ratio.csv}。\file{REF_H} が基準の厚み）。理想的なスリット流れなら 1。 & 倍 & 3e \\
\end{symtab}

% =====================================================================
\section{温度シフト（手順5a）}

\begin{symtab}
$a_T$ & シフトファクター。曲線をどれだけずらすかの係数。基準温度で 1。アレニウス型：$\ln a_T = (E/R)(1/T - 1/T_{\mathrm{ref}})$。 & ― & 5a \\
$E$ & 活性化エネルギー（構造 H・F。横と縦で同じ値）。PP の流動で 40 kJ/mol 前後が文献値の目安。 & kJ/mol & 5a \\
$\Eh$ & 横方向（せん断速度の向き）のシフトの $E$。「見かけのシフトエネルギー（横）」。構造 B・Hz で推定。A では 0。 & kJ/mol & 5a \\
$\Ev$ & 縦方向（粘度の向き）のシフトの $E$。「見かけのシフトエネルギー（縦）」。構造 A・B で推定。Hz では 0。190 ℃ と 230 ℃ の粘度の比は $\exp(0.0207\,\Ev)$（縦だけのとき）。 & kJ/mol & 5a \\
$R$ & 気体定数。 & 8.314 J/mol/K & 5a \\
$T_{\mathrm{ref}}$ & シフトの基準温度（中間の 210 ℃）。ここで $a_T=1$。 & 483.15 K & 5a \\
$f$ & ガウス過程（GP）で表す、基準の流動曲線（マスターカーブ）。 & ― & 5a \\
$u$ & 横シフト後の入力 $u = \ln\gd + \ln a_T(\Eh)$。 & ― & 5a \\
$\ln\mathrm{Fo}$ & $f$ の第2の入力（冷却の度合い。$w$ とも書く）。 & ― & 5a \\
$\varepsilon$ & 観測のノイズ（ショット間のばらつき）。 & ― & 5a \\
マスターカーブ & シフトで換算したあとの、すべての温度のデータを重ねた1本の曲線。重なれば、その構造でシフトできている。 & ― & 5a \\
190/230 ℃ の粘度比 & 同じ厚み・同じ $\gd$ での $\eta(190)/\eta(230)$。モデルでは、Tset だけを 190 ℃・230 ℃ にして予測した粘度の比。\file{ratio_190_230}。実測は 1.07〜1.20 倍で、厚いほど大きい。 & 倍 & 3b・3c・5a \\
LML & 対数周辺尤度。データがそのモデルで、どれだけ「ありそう」か。大きいほどよい。差が 2〜3 以上なら、はっきりした違い。 & ― & 5a（3a〜） \\
$\Delta$LML & 温度なし（$E=0$）との LML の差（\file{dLML_vs_E0}）。正なら、シフトのほうがデータに合う。 & ― & 5a \\
プロファイル尤度 & $E$ を格子で振り、各 $E$ で他を最適にしたときの LML。$E$ と LML の関係の図（尤度曲線）。 & ― & 5a \\
1.92 & LML が最大から 1.92 以内の範囲（\file{E_95_lo}〜\file{E_95_hi}）。おおよそ 95\% 信頼区間。 & ― & 5a \\
粗い格子・細かい格子 & $E$ の格子。粗い格子：1 次元は 2.5 kJ/mol、B は 5 kJ/mol 刻み。細かい格子：最良の点の近くを 0.5（B は 1）kJ/mol 刻み。 & ― & 5a \\
本命・対案 & 5a の決定（2026-10-12）。本命は構造 B、対案は構造 A。基準は構造 C。ノートの設定エリア \file{MAIN_STRUCT}・\file{ALT_STRUCT} で変えられる。 & ― & 5a \\
\end{symtab}

% =====================================================================
\section{構造 × H の扱い（手順5）}

\begin{symtab}
構造 B・A・C & B：横と縦を別々にシフト（本命）。A：縦だけシフト（対案。$\Eh=0$）。C：構造なし（$\Tset$ を入力に並べた通常の GPR）。 & ― & 5 \\
H の扱い (1)・(2)・(4) & ユーザー決定（10-14）。(1)：$H$（Fo）を残す。(2)：材料の部分 $m$ と厚み・冷却の補正 $c(\mathrm{Fo})$ に分ける。(4)：\PII\ を圧力として入れる（$H$ なし）。(3)（$H$ も圧力も外す）は除く。 & ― & 3e・5 \\
(1+4)、(2+4) & (1) と (4) の組み合わせ（$H$ と \PII\ の両方）、(2) と (4) の組み合わせ。ノートの \file{hv} の値は、\file{"1"}・\file{"2"}・\file{"4"}・\file{"14"}・\file{"24"}。 & ― & 5 \\
B\_1、A\_2、C\_4 など & 「構造\_H の扱い」の形のモデル名（\file{model} 列）。手順3b の B0・B2 とは別物。15 個（3 構造 $\times$ 5 通り）。 & ― & 5 \\
$w$ & $w=\ln\mathrm{Fo}=-(\ln\gd+3\ln H)$ ＋ 定数。$H$ と $\gd$ だけで計算できる。 & ― & 5 \\
$p$ & $p=\ln\PII$（条件平均、Pa）。(4) を含むモデルの入力。 & ― & 5 \\
$m$、$c$ & (2) の足し算の成分。$m$：材料の部分（$\gd$・温度〔・$p$〕の関数。Cross-WLF にあたる）。$c$：厚み・冷却の補正（$\ln\mathrm{Fo}$ の関数）。$m$ と $c$ は定数だけ入れ替えられるので、図では $c$ の平均を 0 にそろえる。 & ― & 5 \\
\file{OnDims} & 入力の一部の列だけに当てるカーネル（足し算のカーネル $m+c$ を作るための部品。\file{nb_common.py}）。 & ― & 5 \\
実測の \PII、予測した \PII & (4) を含むモデルの評価の 2 通り。実測：検証の条件の実測の条件平均を使う（\PII\ が分かっている場合の上限）。予測：$(\ln\gd,\ln H,\Tset)$ から GPR で予測し、予測分布から \file{P32_DRAWS} 回引いて平均と分散にまとめる（\file{psrc} 列の「実測」「予測」）。 & ― & 5 \\
$E_v-E_h$（\file{diff}） & B で決まる量。$\Eh$ と $\Ev$ は別々には決まりにくく（尤度の山が斜めの帯）、この差だけが決まる。 & kJ/mol & 5a・5 \\
\end{symtab}

% =====================================================================
\section{最終候補の MCMC（手順5b）}

\begin{symtab}
MLE（E 固定）・MLE（同時） & 手順5b の最尤推定。E 固定：ずらし量を手順5の最良値に固定してカーネルだけ推定（手順5と同じ設定）。同時：ずらし量もカーネルと同時に推定（MCMC の初期値）。\file{mle.csv}・\file{metrics.csv} の \file{method}。 & ― & 5b \\
条件ごとにまとめた尤度 & 入力が条件平均なので、ショットごとの尤度を、条件ごとの平均・ショット数・条件内の平方和で厳密に書き直したもの（\file{step5b_core.py} の \file{grouped_lml}）。条件の数の大きさの行列だけで済み、1 回約 1 ミリ秒。ショットごとの GPR と一致する。 & ― & 5b \\
$\varepsilon$、$\sigma_n^2$ & ノイズ（ショット間のばらつき）と、その分散（WhiteKernel。\file{noise}・\file{log_noise}）。 & ― & 5・5b \\
事前分布（E） & MCMC での $\Eh$・$\Ev$ の事前分布。範囲の中で一様（\file{E_PRIOR}：$\Eh$ $-40$〜$40$、$\Ev$ $-20$〜$40$ kJ/mol）。 & kJ/mol & 5b \\
\file{amp}、\file{lin_amp}、\file{ls_<入力>}、\file{sigma0}、\file{noise} & posterior\_summary.csv の \file{param}（元の単位。振幅は $\sqrt{\text{振幅}^2}$）。入力の記号は \file{u}・\file{w}・\file{H}・\file{T}・\file{p}。 & ― & 5b \\
\file{MLE_E_fixed}、\file{MLE_joint}、\file{profile_best_E} & posterior\_summary.csv：最尤推定（E 固定／同時）の値、手順5のプロファイル尤度の最良の E。 & ― & 5b \\
\file{method} & metrics.csv・tset\_effect.csv：予測の方法（MLE（E 固定）、MLE（同時）、MCMC）。MCMC は、事後サンプルの予測分布の混合分布。 & ― & 5b \\
\file{MCMC_MODELS}、\file{SETS5B} & 手順5b の設定：MCMC で推定するモデル（\file{["B_1", "C_1"]}）と、版ごとの学習データ。 & ― & 5b \\
\file{N_WARM}、\file{N_SAMP}、\file{N_SAMP_MAX}、\file{N_PRED_DRAWS}、\file{E_START} & 手順5b の設定：調整の反復数（5000）、本番の反復数（20000）、本番の上限（100000）、予測に使う事後サンプル数（200）、手順5の結果がないときの E の出発点。 & ― & 5b \\
\end{symtab}

% =====================================================================
\section{材料別 GPR（05 フォルダ。本体とは別のプロジェクト）}

\file{05_材料ごとの機械学習データと予測検証データ} のデータだけを使う、別のプロジェクト（2026-10-14）。説明変数に圧力 P3\_3 と $\Tff$ を使う点が、本体（CLAUDE.md §5）と違う（ユーザー決定）。説明書は \file{05_材料ごとの機械学習データと予測検証データ/学習結果説明_材料別温度シフトGPR.pdf}。「出る所」の「材」は、この説明書。

\begin{symtab}
P3\_3 & 流動中樹脂圧力 P3\_3（列 \file{流動中樹脂圧力P3_3(MPa)}、ノートの \file{P33}・\file{ln_P33}〔Pa に直して ln〕）。05 のデータで増えた列。粘度の算出式の圧力ではない（算出式の形で表すと誤差が数十〜数百 \%）。ln P3\_2 との相関は 0.95〜1.00。 & MPa & 材 \\
GPPS & 材料（非晶性）。05 の \file{03_GPPS実験}。設定温度 200・220・240 ℃、厚み 1.0・1.5・2.0 mm。P3\_1・P3\_2 の列はない。 & ― & 材 \\
B\_Tset、B\_Tff & 材料別 GPR のシフトのあるモデル。B\_Tset：$\ln\eta=\ln a_v(\Tset)+f(\ln\gd+\ln a_h(\Tset),\ln\mathrm{P3\_3},\ln\Tff)$。B\_Tff：$\ln\eta=\ln a_v(\Tff)+f(\ln\gd+\ln a_h(\Tff),\ln\mathrm{P3\_3},\Tset)$。本体の B\_1 とは別物。 & ― & 材 \\
A\_Tset、A\_Tff & 縦だけシフトする対案（$\theta_h=0$）。設定 \file{STRUCTS = ["B", "A"]} で作る。 & ― & 材 \\
C（材料別） & シフトの構造なしで、$\ln\gd$・$\ln$ P3\_3・$\ln\Tff$・$\Tset$ の 4 つを並べた通常の GPR（6月のノートの形に $\Tset$ を足したもの）。本体の構造 C とは入力が違う。 & ― & 材 \\
$a_h(T)$、$a_v(T)$ & 横（せん断速度の向き）と縦（粘度の向き）のシフトファクター。$T_{\mathrm{ref}}$ で 1。 & ― & 材 \\
$\theta_h$、$\theta_v$（\file{th}・\file{tv}） & シフトファクターの係数。アレニウス型では $\Eh$・$\Ev$（kJ/mol）、WLF 型では $A_{1,h}$・$A_{1,v}$。CSV の \file{th_best}・\file{tv_best}・\file{th_lo}〜\file{tv_hi}（95\% 区間の目安）・\file{diff_best}（$\theta_v-\theta_h$）・\file{dLML_vs_0}（$\theta=0$ との LML の差）。 & kJ/mol／― & 材 \\
アレニウス型、WLF 型 & シフトファクターの式。アレニウス型：$\ln a=(E/R)(1/T-1/T_{\mathrm{ref}})$（PP・粉砕材混合 PP）。WLF 型：$\ln a=-A_1[(T-D_2)/(A_2+T-D_2)-(T_{\mathrm{ref}}-D_2)/(A_2+T_{\mathrm{ref}}-D_2)]$（GPPS）。 & ― & 材 \\
$A_1$、$A_2$、$D_2$、$T_g$ & WLF 型の係数。$A_2=51.6$ K と $D_2=T_g$（ガラス転移温度。仮置き 100 ℃）は固定し、$A_1$ を推定する。参考の Cross-WLF では、$D_2$ は K で、PP は $A_1=30.9$・$A_2=43.4486$ を固定（6月のノートと同じ）。 & ―／K／℃ & 材 \\
$T_{\mathrm{ref}}$（材料別） & 基準温度。PP・粉砕材混合 PP は 210 ℃、GPPS は 220 ℃（中央の設定温度）。 & ℃ & 材 \\
規則 g & P3\_3 が、同じ条件の中央値から 20\% 以上離れたショット（問題の候補。\file{P33_DEV}）。 & ― & 材 \\
尾根、差 $\theta_v-\theta_h$（\file{diff_best}・\file{diff_lo}・\file{diff_hi}） & プロファイル尤度の山が、斜めの帯（尾根）になること。$\theta_h$ と $\theta_v$ は別々には決まりにくく、差 $\theta_v-\theta_h$ が最もよく決まる（本体の手順5a・5 と同じ）。個々の係数の 95\% 区間は、計算した範囲に左右される。 & ― & 材 \\
尤度比検定の 5\% 基準（\file{LRT_DLML}） & 2 つの係数（横・縦のシフト）を足した効果が認められる目安。$2\times\Delta\mathrm{LML}>5.99$（自由度 2 のカイ二乗分布の 95\% 点）、つまり $\theta=0$ との LML の差が 3.0 以上。ノートの判定に使う。 & ― & 材 \\
\file{refine_converged}、\file{MAX_REFINE} & 細かい格子の取り直し：最良の点が動いたら、その点を中心に細かい格子を取り直す。収束（最良の点が動かない）したか、取り直しの回数の上限（既定 4）。 & ― & 材 \\
Tboth（Tset と Tff を同時に） & 依存性グラフの図と判定：設定温度 $\Tset$ を上げると、流動中樹脂温度 $\Tff$ も同じだけ上がるとして、2 つを一緒に動かす（$\Tff=\Tff^{(0)}+(\Tset-\Tset^{(0)})$）。成形で起こることに近い。$\Tset$ だけ・$\Tff$ だけを動かす図は、ほかの変数の固定の仕方によって、効果が小さく・逆向きに出ることがある。 & ― & 材 \\
念入りな最適化（5.1b）、\file{N_RESTARTS} & 最良の $\theta$ と $\theta=0$ の 2 点だけを、最適化のやり直し \file{N_RESTARTS}（既定 20。6月は 5）回で学習し直すこと。格子の各点の LML は、ハイパーパラメータの局所解で 1〜3 ずれるため。最良の $\theta$ の学習が、そのまま最終モデル。\file{dLML_vs_0} は、この 2 つの差（格子での差は \file{dLML_vs_0_grid}）。 & ― & 材 \\
\file{mat_seed_check.py}、\file{乱数シードの確認.csv} & 補助スクリプト（ノートの外）：乱数シードを変えて最終モデルを学習し直し、LML と 3data の RMSE のばらつきを見る。結果は \file{<入力>/比較/乱数シードの確認.csv}・\file{…_まとめ.csv}・\file{乱数シードの確認.png}。 & ― & 材 \\
\file{FLOW_COMMON}、\file{FLOW_COMMON_XLIM}・\file{YLIM}・\file{NCOL} & 全厚み共通のせん断速度 ― 粘度の図の設定。温度ごとの図 \file{…_全厚み共通_<温度>C.png} と、全温度を 1 枚に重ねた総合図 \file{…_全厚み共通_全温度.png}（線・帯の色は温度：赤・青・緑。グラフに凡例・タイトルはなく、凡例は \file{…_凡例.png}）。6月のセル9-2 と同じ仕様（15 $\times$ 10 インチ、10 の累乗の目盛り）。 & ― & 材 \\
\file{CWLF_STARTS} & 参考の Cross-WLF の当てはめの初期値の数（6月に近い値 1 つ ＋ 乱数。残差が最小のものを採る）。 & ― & 材 \\
重ね合わせの残差 & 厚みごとに、全温度の条件平均の点を 3 次式で当てはめた残差の RMSE（ln η）。シフトで温度ごとの曲線が重なるほど小さい。 & ― & 材 \\
P09 & 05 の問題データ一覧：GPPS の 3data、200 ℃・1.5 mm・3 mm/s のショット 2・3（P3\_3 が負）。「ショット」列に \file{2;3} と複数書ける（材料別 GPR のノートだけが読む書き方）。 & ― & 材 \\
\file{MATERIALS}、\file{SHIFT_TEMPS}、\file{STRUCTS}、\file{MAIN_KEY} & 材料別 GPR の設定：学習する材料、シフトに使う温度（\file{"Tset"}・\file{"Tff"}）、シフトの形（\file{"B"}・\file{"A"}）、図で主に見せるモデル。 & ― & 材 \\
\file{SHIFT_CFG}、\file{FINE_B}、\file{PROFILE_RESTARTS} & 材料ごとのシフトファクターの型・$T_{\mathrm{ref}}$・$T_g$・尤度の格子（粗い格子の範囲と刻み）、細かい格子、格子の各点でのやり直し回数。 & ― & 材 \\
\file{OUTLIER_REGISTER}、\file{CWLF_T}、\file{CWLF_FIXED}、\file{USE_LINEAR} & 外れ値を一覧に登録して除く材料（PP・MIX）、Cross-WLF に使う温度と固定する係数、線形項を足すか。 & ― & 材 \\
\file{INPUT_MODE}（ショットごとの入力・条件平均の入力） & $\gd$・P3\_3・$\Tff$ の入れ方。\file{"shot"}（\textbf{既定}。ユーザー決定、2026-10-23）：ショットごとの値（6月のノートと同じ入れ方。1 ショットを 1 つのデータとして学習）。\file{"cond"}（参考。10-15〜10-22 の既定）：条件ごとの平均（列 \file{ln_gd_cm}・\file{ln_P33_cm}・\file{ln_Tff_cm}・\file{Tff_cm}。本体の手順3b の R3 と同じ考え方）。結果は \file{<材料フォルダ>_学習結果/ショットごとの入力/}・\file{条件平均の入力/} に分けて保存する。 & ― & 材 \\
\file{REPO}（リポジトリ直下） & ノートの設定エリアで、ノートの場所から上へたどって、\file{04_補助スクリプト} があるフォルダを自動で探す（2026-10-23。ノートは \file{01_機械学習コードフォルダ/材料別_温度シフトファクター検討/} に移した）。 & ― & 材 \\
\file{<材料フォルダ>_学習結果/} & 結果の置き場所（学習データ・検証データのフォルダと同じ場所）。中に \file{データ確認/}・\file{B_Tset/}・\file{B_Tff/}・\file{C/}・\file{比較/}。600 dpi の図は各フォルダの \file{highres/}（Git の管理外）。 & ― & 材 \\
厚みごとの比較（ノート） & 材料別 GPR の結果を、キャビティ厚みごとに調べ直す別のノート \file{材料別_温度シフトGPR_ARD対応Matern52＋線形＋ホワイトノイズ_厚みごとの比較.ipynb}（2026-10-22）。元のノートの学習済みモデルと予測値を読む。結果は \file{<材料フォルダ>_学習結果/<入力の取り方>/厚みごとの比較/}（既定はショットごとの入力）。部品は \file{04_補助スクリプト/mat_thickness.py}、生成は \file{build_mat_thick_nb.py}。 & ― & 厚 \\
金型内樹脂最高温度 $T$ & 発表資料（報告書B・ポスター）での $\Tff$（流動中樹脂温度）の呼び名。同じもの。 & ℃ & 厚 \\
全厚み共通の係数、厚みごとの係数 & 全厚みの 5data で 1 組だけ決めたシフトの係数（元のノートのモデル）と、厚みごとに決め直した係数。 & kJ/mol／― & 厚 \\
$R$（温度による粘度の比。列 \file{R_実測}・\file{R_予測}・\file{lnR の差}・「比べた $\gd$ の範囲」） & 同じせん断速度での、最低と最高の加熱シリンダ温度の粘度の比。成形条件ごとの平均を、$\ln\gd$ の重なる範囲の共通の格子（$G=50$ 点）で内挿し、$R=\exp(\frac{1}{G}\sum_g[\ln\eta_{\mathrm{low}}(\gd_g)-\ln\eta_{\mathrm{high}}(\gd_g)])$。予測の $R$ は、各条件の予測平均から同じ手順で出す。本体のモデル名 R1〜R3 や、気体定数 $R$ とは別物。 & ― & 厚 \\
重ね合わせのずれ（厚みごと。列 \file{ずれ_換算前}・\file{ずれ_換算後}・\texttt{ずれ\_Tset\_\%}・\texttt{ずれ\_Tff\_\%}） & 1 つの厚みの条件平均の点を、係数で換算し、換算した $\ln\gd$ に対する 3 次式を当てはめた残差の RMSE（$\ln\eta$）。\texttt{\_\%} の列は粘度の \% に直したもの（$(e^{s}-1)\times100$）。重ね合わせの残差を、厚みごとに分けたもの。 & ―／\% & 厚 \\
判定 ○・△・×（列 \file{順序が一致}・\file{比の大きさが一致}・\file{換算でずれが減る}・\file{判定}） & 厚みごとの温度シフトの確認。○：温度の順序が一致し、$|\ln R_{\text{予測}}-\ln R_{\text{実測}}|$ が許容差以内で、換算でずれが減る。×：予測の順序が実測と逆。△：それ以外。 & ― & 厚 \\
許容差（列「許容差（ショット間の標準偏差）」） & 同じ成形条件のショット間の $\ln\eta$ の標準偏差（純粋誤差。条件ごとの分散を自由度で重み付けしてまとめたもの。約 0.04）。2026-10-22 の版はモデルのノイズの標準偏差を使ったが、ショットごとの入力ではノイズが約 0.006 と小さく出るので、入力の取り方によらないこの値に変えた（10-23）。列 \file{ノイズの標準偏差} は参考。 & ― & 厚 \\
(a)・(b)・(c)、列 \file{適した温度(a)}〜\file{(c)}・\file{総合の判定} & シフトに使う温度を厚みごとに比べる 3 つの基準：(a) 重ね合わせ（厚みごとの係数で最小のずれ）、(b) 厚みごとの 5data での LML の差 $\Delta L$、(c) 3data の予測誤差。値は「Tff」「Tset」「差なし」「—」。総合の判定は 3 つの多数決（同数なら差なし）。 & ― & 厚 \\
$\Delta L$（列「$\Delta$L（Tff − Tset）」） & $\mathrm{LML}(\text{B\_Tff})-\mathrm{LML}(\text{B\_Tset})$。同じ厚みの 5data で、それぞれの係数をプロファイル尤度で決めたときの最大の LML の差。$|\Delta L|\ge3$ で差あり。 & ― & 厚 \\
\file{係数0との差}（\file{プロファイル尤度_厚みごと.csv}） & 厚みごとの 5data で、最良の係数と係数 0（シフトなし）の LML の差。元のノートの \file{dLML_vs_0} の厚みごと版（念入りな最適化はしない）。 & ― & 厚 \\
\file{A_CLEAR}、\file{DL_CLEAR}、\file{ERR_CLEAR} & (a)(b)(c) で「差あり」とする差の大きさ：ずれ 0.5 ポイント、$|\Delta L|$ 3.0、予測誤差 1.0 ポイント（既定）。 & \%／―／\% & 厚 \\
\file{THICK_PROFILE}、\file{THICK_MODELS} & 厚みごとのプロファイル尤度 (b) を計算するか、厚みごとのモデル（9 節）を学習するか。 & ― & 厚 \\
厚みごとのモデル、総合データのモデル（列 \file{学習} が \file{厚みごと}・\file{総合}） & 1 つの厚みの 5data だけで学習したモデル（係数は (b) で決めたもの）と、全厚みの 5data で学習した元のモデル。同じ 3data で比べる。列 \file{厚みの名前} は厚み（\file{0.5}・\file{1} など）か \file{全厚み}、列 \file{範囲} は \file{全体}・\file{低速側}・\file{高速側}。 & ― & 厚 \\
\texttt{誤差\_\%}（予測誤差） & 3data の $\ln\eta$ の RMSE を、粘度の \% に直した値 $(e^{\mathrm{RMSE}}-1)\times100$。 & \% & 厚 \\
厚みごとの比較の CSV & \file{温度シフトの確認_厚みごと.csv}（問い 1）、\file{重ね合わせ_厚みごと.csv}（(a)）、\file{プロファイル尤度_厚みごと.csv}・\file{…_格子_厚みごと.csv}・\file{対数周辺尤度の差_厚みごと.csv}（(b)）、\file{検証誤差_厚みごと.csv}（(c)）、\file{厚みごとの比較_まとめ.csv}、\file{厚みごとと総合の比較.csv}・\file{…_まとめ.csv}・\file{厚みごとのモデル_予測値.csv}（問い 3）、\file{判定の一覧.csv}。 & ― & 厚 \\
\end{symtab}

% =====================================================================
\section{データの呼び名}

\begin{symtab}
PP、MIX & 材料。PP はバージン PP。MIX（粉砕材混合 PP）は、PP に一次成形の PP（粉砕材）を 25 wt\% 加えたもの。 & ― & 3a〜 \\
5data & 学習用のファイル（各条件 4〜5 ショット。93 条件）。 & ― & 3a〜 \\
3data & 検証用のファイル（各条件 1〜3 ショット。36 条件）。5data と条件が重ならず、射出速度の両端にある。 & ― & 3a〜 \\
Alldata & 学習と検証をまとめたファイル（PP は各条件 10 ショット。129 条件）。従来のランダム分割の評価に使った（手順3a・3b）。手順3d 以降は、「Alldata 版」の学習データ。 & ― & 3a〜 \\
外れ値なし・あり & 外れ値（Tff は修正 $z$ が 3.5 超、ほかの列は中央値から 20\% 以上離れた値）を除いたファイルと、除かないファイル。主には「なし」を使う。\textbf{手順3d 以降は、外れ値ありのデータは使わない}（ユーザー決定、2026-10-14）。 & ― & 3a〜 \\
5data 版・Alldata 版 & 手順3d 以降のデータの版（上の「紛らわしい記号」を参照）。どちらも外れ値なしで、問題データ一覧の「除外」を除く。ノートの設定 \file{DATA_VERSIONS} で選ぶ。 & ― & 3d〜 \\
旧 00 & 旧データ（\file{学習データ6月発表時.csv}）。手順2だけで使った。今後は使わない。 & ― & 2 \\
no1.5-3.0 & PP の厚み 1.5 mm・射出速度 3 mm/s の3条件を除いたファイル。ノートは、全条件のファイルを読み、問題データ一覧の「除外」で同じに除く。 & ― & 3a〜 \\
no0.5-320-190 & PP の 190 ℃・0.5 mm・320 mm/s を除いたファイル（手順3b での感度確認の記録。P03 は 10-14 に「除外」になったので、ノートは全条件のファイルから、一覧で除く。このファイルは使わない）。 & ― & 3b \\
P01〜P08 & 問題データ一覧の ID。P01：PP 1.5 mm・3 mm/s（除外）、P02：0.5 mm・20 mm/s・190 ℃（感度確認）、P03：PP の 190 ℃・0.5 mm・320 mm/s（手順3b では感度確認、\textbf{10-14 に除外へ変更}。規則 f で拾った）、P04：Tff が旧データと同じ（注意）、P05：旧 00（使わない）、P06：外れ値の判定の単位が未確認（注意）、P07：外れ値なしは、ばらつきが小さめ（注意）、P08：粉砕材混合 PP の 004\_190 ℃ の行のコピー（除外）。 & ― & 3a〜 \\
除外・感度確認・注意・使わない & 問題データ一覧の区分。除外：解析から除く。感度確認：主の分析に残し、除いた場合を並べる。注意：結果を読むとき気をつける。使わない：今後の解析に使わない。 & ― & 3a〜 \\
再測定 & 一覧の列。未定・予定・済み。「済み」にした行は、除外・感度確認でも使われない（取り直したデータを使うため）。 & ― & 3a〜 \\
\file{is_sens}、\file{sens} & ノートの列。「感度確認」の条件の行に印を付ける。\file{sens} は、その ID（P02 など）。 & ― & 3b〜 \\
結合ファイル & 圧力と温度のデータを結合したファイル（\file{結合_<番号>_<材料>_max_<組>_外れ値<あり|なし>.csv}）。 & ― & 3a〜 \\
\end{symtab}

% =====================================================================
\section{データの分け方}

\begin{symtab}
S1 & 手順2。せん断速度 500〜1000 1/s の 90 行（9 条件）をテストにし、残りで学習。現行ノートと同じ。 & ― & 2 \\
S2 & 手順2。ランダムに 2 割（258 行）をテストにする。評価が甘くなる分け方。 & ― & 2 \\
S\_V & 5data で学習し、3data で検証する（射出速度の両端への外挿）。低速側（3・5 mm/s）と高速側（280・320 mm/s）に分けて評価する。 & ― & 3a〜 \\
S\_T & 5data の 190・230 ℃ で学習し、5data の 210 ℃ の全条件をテストにする（設定温度の内挿）。 & ― & 3b〜 \\
S\_T3 & S\_T と同じ学習で、3data の 210 ℃ をテストにする（温度の内挿に、速度の外挿が重なる）。 & ― & 3b〜 \\
S\_R1、S\_R2 & 従来の、行単位のランダム 20\%。S\_R1：Alldata、S\_R2：PP の 5data＋3data（PP のみ）。\file{random_state}=42。 & ― & 3a〜 \\
LOCO & 条件を1つずつ抜く交差検証（Leave-One-Condition-Out）。ハイパーパラメータは全データで決めて固定し、閉じた式で計算（近似）。 & ― & 3a・3b \\
ALL & 5data＋3data（129 条件）。5a で、最終的な $E$ の推定値を見るために使う。手順3d 以降は、「最終モデル」の学習データ（5data 版：5data、Alldata 版：Alldata の全条件）。 & ― & 5a／3d \\
S\_Vall & Alldata 版の外挿の検証。Alldata のうち 5data の条件（全ショット）で学習し、3data を検証する（PP だけ意味がある。粉砕材混合 PP は、5data と同じデータなので S\_V と同じ）。 & ― & 3d \\
CV（\file{CV1}〜\file{CV10}）& 条件単位の 10 分割の交差検証。同じ条件のショットは、同じ側に入る。各分割の予測をまとめて、全体と、低速側・内側・高速側で指標を出す。Alldata 版と、5data 版（5data の中）で行う。 & ― & 3d \\
S\_T（Alldata 版）& Alldata の 190・230 ℃ で学習。テストは 5data 版と同じ行（5data ファイルの 210 ℃ と、3data ファイルの 210 ℃ = S\_T3）なので、そのまま比べられる。 & ― & 3d \\
低速側・高速側・内側 & 検証データの区分。低速側：射出速度 3・5 mm/s、高速側：280・320 mm/s、内側：その間。 & ― & 3a〜 \\
感度除く・（感度確認の条件を除く） & 結果名の後ろの印。学習から（またはテストの評価から）「感度確認」の条件を除いた場合。 & ― & 3b・5a \\
cm、shot & 入力の $\gd$ の入れ方。cm：条件平均（成形条件から見た予測）、shot：ショットごと（実測の $\gd$）。 & ― & 3b〜5a \\
\end{symtab}

% =====================================================================
\section{モデル名}

\subsection{名前の付け方（手順3b 以降）}
\begin{points}
\pt{先頭の B と R}B 系：ショットごとの $\gd$（\file{ln_gd}）を入力に使う GPR。R 系：条件ごとの平均の $\gd$（\file{ln_gd_cm}）を使う GPR（Reference：設定値から見た形）。ただし R2 は、射出速度 $v$ を使う。
\pt{末尾の T・F・P}T：$\Tset$ を入力に足した。F：$\Tff$ を足した。P：\PII\ を足した（R3P。B2 系の 2 は \PII）。FT は、$\Tff$ と $\Tset$ の両方。
\pt{L}線形回帰。L0：$\gd$・$H$ だけ。L0T：さらに $\Tset$。
\end{points}

\begin{symtab}
M1 & 現行の GPR。入力 $\ln T$・$\ln P$・$\ln\gd$（手順2）。 & ― & 2 \\
M2 & 線形回帰。入力 $\ln P$・$\ln\gd$。 & ― & 2 \\
M3 & 線形回帰（算出式どおり）。入力 $\ln(P-P_2)$・$\ln\gd$・$\ln H$。 & ― & 2 \\
M4 & $P$ を外した GPR。入力 $\ln T$・$\ln\gd$。 & ― & 2 \\
L0・L1・L2 & 線形回帰。入力：$\gd$・$H$ ／ ＋\PI ／ ＋\PII（手順3a）。 & ― & 3a \\
A0・A1・A2 & GPR。入力：$\gd$・$H$・$\Tff$ ／ ＋\PI ／ ＋\PII（手順3a）。 & ― & 3a \\
A1$'$ & GPR。入力：$\gd$・$\Tff$・\PI（現行ノートと同じ形。厚み $H$ がない）。 & ― & 3a \\
L0T & 線形回帰。入力：$\gd$・$H$・$\Tset$。 & ― & 3b・3c \\
B0・B0T & GPR（ショットごと）。入力：$\gd$・$H$ ／ ＋$\Tset$。 & ― & 3b・3c \\
B0F・B0FT & GPR（ショットごと、$\Tff$ あり）。入力：$\gd$・$H$・$\Tff$ ／ ＋$\Tset$。 & ― & 3b \\
B2・B2T & GPR（ショットごと、\PII\ あり）。入力：$\gd$・$H$・\PII ／ ＋$\Tset$。 & ― & 3b \\
B2F・B2FT & GPR（ショットごと、$\Tff$・\PII\ あり）。 & ― & 3b \\
R0 & GPR（条件ごと）。入力：条件平均の $\gd$・$H$。 & ― & 3b \\
R2 & GPR（設定値だけ）。入力：$v$・$H$・$\Tset$。成形条件だけで予測できるか。 & ― & 3b・3c \\
R3 & GPR（条件ごと）。入力：条件平均の $\gd$・$H$・$\Tset$。基本のモデル。CLAUDE.md §6 の R1。 & ― & 3b〜5a \\
R3F・R3P & R3 に、$\Tff$（条件平均）／\PII（条件平均）を足した GPR。 & ― & 3b \\
H・A・Hz・B & 温度シフトの構造（手順5a）。H：横も縦も同じ $E$（当初の構造）。A：縦だけ。Hz：横だけ。B：横と縦を別々。 & ― & 5a \\
N・H40・F・C & N：温度なし（$E=0$）。H40：H で $E=40$ kJ/mol に固定。F：$\Tff$ でシフト（参考）。C：構造なしの通常の GPR（$\gd$・$H$・$\Tset$ を並べる）。 & ― & 5a \\
RFo & R3 の $H$ を、フーリエ数に替えた GPR。入力：条件平均の $\gd$・$\ln\mathrm{Fo}$・$\Tset$（\file{ln_Fo_cm}）。$H$ と同じ情報。 & ― & 3e \\
W0 & Cross-WLF と同じ変数の GPR（温度は設定温度）。入力：条件平均の $\gd$・$\Tset$。$H$ を外したモデル。W は WLF の頭文字。 & ― & 3e \\
W1・W1T & Cross-WLF と同じ変数の GPR（温度は流動先端温度）。W1：条件平均の $\gd$・$\ln\Tff$。W1T：さらに $\Tset$。 & ― & 3e \\
W2・W3 & Cross-WLF と同じ変数の GPR（圧力 $P$ として \PII\ を入れる。$H$ なし）。W2：条件平均の $\gd$・$\Tset$・$\ln\PII$。W3：条件平均の $\gd$・$\ln\Tff$・$\ln\PII$。精度は R3 に近いが、\PII\ が厚み（$H^{-1.5}$ に比例）を運んでいるため（手順3e）。 & ― & 3e \\
B\_1〜C\_24 & 手順5 のモデル（構造\_H の扱い。上の「構造 $\times$ H の扱い（手順5）」）。C\_1 は R3、C\_4 は W2 と同じ。 & ― & 5 \\
\end{symtab}

% =====================================================================
\section{評価・推定・MCMC}

\begin{symtab}
$R^2$ & 決定係数。1 なら完全に一致、0 なら「いつも平均を答える」のと同じ、マイナスならそれより悪い。 & ― & 全 \\
RMSE & 二乗平均平方根誤差。予測と実測の差の平均的な大きさ（ここでは $\ln\eta$）。0.05 なら粘度で約 5\%。 & ― & 全 \\
被覆率（\file{coverage95}）& 95\% 予測区間に実測が入った割合。0.95 が理想。 & ― & 全 \\
NLPD & 負の対数予測密度。当たり具合と区間の適切さをまとめた指標。小さいほどよい。 & ― & 全 \\
\file{mean_sd} & 予測標準偏差の平均（区間の広さの目安）。\file{mean_sd_latent} は、ノイズを除いた（潜在関数の）もの。 & ― & 3a〜 \\
\file{noise_sd} & GPR のノイズの標準偏差（$\sqrt{\text{ノイズ分散}}$）。同じ条件でのショット間のばらつきを表す。ショットごとの $\gd$ を使うと極端に小さくなる（$5\times10^{-5}$ ほど）。 & ― & 3a〜 \\
予測区間・信頼区間 & 予測区間：新しく測った値が入る範囲（ノイズ込み）。信頼区間：平均の不確かさだけ（ノイズなし）。この資料では、図の帯は 95\% の予測区間。 & ― & 全 \\
標準化・中心化 & 標準化：平均を引き、標準偏差で割る（入力）。中心化：平均を引く（$\ln\eta$）。学習データだけで計算する。 & ― & 全 \\
内挿・外挿 & 学習範囲の内側／外側の予測。 & ― & 全 \\
リーク & 答えの計算に使った値（\PI）が入力に紛れ込み、答えを「計算」できてしまうこと。 & ― & 2・3a \\
MCMC & マルコフ連鎖モンテカルロ法。ハイパーパラメータの事後分布を、サンプルで表す。自作の適応型メトロポリス法。 & ― & 2・3c \\
$\theta$ & ハイパーパラメータの対数（ln）をまとめたもの。\file{log_c}・\file{log_ls_<入力>}・\file{log_c_lin}・\file{log_sigma0}・\file{log_noise}。 & ― & 3c \\
$\lambda$、$\Sigma$ & 適応型メトロポリス法の提案分布 $N(\theta,\lambda\Sigma)$ の、大きさ（$\lambda$）と形・向き（$\Sigma$）。 & ― & 3c \\
採択率 & 提案に移れた割合。目標 0.234。 & ― & 2・3c \\
調整期間・本番 & 調整期間（3,000 回）：提案分布を合わせる期間で、捨てる。本番（8,000 回、足りなければ 20,000 回）：結果に使う。 & ― & 2・3c \\
チェーン & 別々の出発点から走らせる一続きのサンプル（4 本）。 & ― & 2・3c \\
R-hat & チェーンがそろっているかの指標。1.01 以下が目標。 & ― & 2・3c \\
ESS & 有効サンプル数。bulk は分布の中心、tail は裾。400 以上が目標。 & ― & 2・3c \\
事前・事後分布 & 事前：データを見る前の分布（探索範囲の中で、対数について一様）。事後：データを見たあと。 & ― & 2・3c \\
混合分布 & 事後サンプルごとの予測分布を、同じ重みで混ぜた分布。ハイパーパラメータの不確かさも、予測の幅に入る。 & ― & 2・3c \\
\end{symtab}

% =====================================================================
\section{カーネルとハイパーパラメータ}

\begin{symtab}
GPR & ガウス過程回帰。似た入力には似た出力が対応する、という考えでデータの間をつなぐ回帰。予測値と一緒に、不確かさも出せる。scikit-learn の \file{GaussianProcessRegressor}。 & ― & 全 \\
Matern 5/2 & ほどよく滑らかなカーネル。入力ごとの長さスケールを持つ（ARD）。 & ― & 全 \\
ARD & 入力ごとに別々の長さスケールを持たせる仕組み。長さスケールを見ると、どの入力がよく効くか分かる。 & ― & 全 \\
長さスケール（\file{ls_<入力>}）& 入力がどれだけ変わると出力が大きく変わるかの目安（標準化した入力に対する値）。小さいほどよく効く。探索範囲は 0.01〜100。 & ― & 全 \\
感度 & 長さスケールの逆数。ARD の図（3b・3a）で使う。 & ― & 3a・3b \\
振幅（\file{amplitude}）& カーネル C の $\sqrt{\text{振幅}^2}$。関数の大きさの目安。探索範囲は 振幅$^2$：$10^{-3}$〜$10^{3}$。 & ― & 全 \\
線形カーネル & DotProduct（$\sigma_0^2 + x\cdot x'$）に振幅の 2 乗を掛けたもの。低速側への外挿のために足す（ユーザー決定、10-12）。\file{lin_amplitude}・\file{lin_sigma0}。 & ― & 3b〜 \\
$\sigma_0$ & 線形カーネルの定数項（\file{lin_sigma0}）。データでほとんど決まらず、探索範囲いっぱいに広がる（予測にはほとんど効かない）。 & ― & 3c \\
ノイズ分散 & WhiteKernel のパラメータ（\file{noise_var}）。探索範囲の下限は \file{NOISE_LB} = $10^{-6}$（旧ノートは $10^{-3}$）。 & ― & 全 \\
張り付き & ハイパーパラメータが探索範囲の端に止まること（\file{noise_at_lower_bound}・\file{ls_at_bound}）。長さスケールが下限に張り付くと、その入力を内挿できない。 & ― & 3a〜 \\
\end{symtab}

% =====================================================================
\section{ノートの設定・列名・ファイル}

\subsection{ノートの設定エリア}
\begin{symtab}
\file{SEED} & 乱数シード（42）。 & ― & 全 \\
\file{NOISE_LB} & ノイズ分散の下限。 & $10^{-6}$ & 3a〜 \\
\file{N_RESTARTS} & GPR の最適化のやり直し回数。 & 5 & 3a〜 \\
\file{USE_EXCLUSIONS} & 問題データ一覧の「除外」を使うか。 & True & 3a〜 \\
\file{MAIN_DATA} & 主に使う外れ値の扱い（\file{なし}）。 & ― & 3b〜 \\
\file{TEMP_HOLDOUT} & S\_T でテストにする設定温度（210）。 & ℃ & 3b〜 \\
\file{MAIN_STRUCT}、\file{ALT_STRUCT} & 5a の本命（B）と対案（A）。 & ― & 5a \\
\file{DATA_VERSIONS} & 検討するデータの版（\file{["5data", "Alldata"]}。片方だけにもできる）。 & ― & 3d〜 \\
\file{CV_K}、\file{N_RESTARTS_CV} & 交差検証の分割数（10）、各分割での最適化のやり直し回数（1）。 & ― & 3d \\
\file{N_JOBS} & 並列に計算する数（CPU の数に合わせる）。 & 6 & 3d・5a \\
\file{DRIFT_TOL} & ショットの順番によるずれの許容値（$\ln\eta$）。 & 0.02 & 3d \\
\file{OUTLIER_Z}、\file{OUTLIER_RESID} & 規則 f（外れ値の検知）のしきい値：$|z|$（3.5）と |残差 $\ln\eta$|（0.1）。 & ― & 3d \\
\file{HR_ALL}、\file{HR_ALL_DPI} & すべての図を、高解像度でも \file{highres/} に保存するか（True）、その解像度（600 dpi）。 & ― & 3d〜 \\
\file{HR_VERSIONS} & 高解像度のグラフ（せん断速度 ― 粘度）で描く版。 & ― & 3d・3e \\
\file{KEYS}、\file{REF_H} & 手順3e で比べるモデル（R3・RFo・W0・W1・W1T・W2・W3）、厚みの比の基準の厚み（2.0 mm）。 & ― & 3e \\
\file{STRUCTS}、\file{H_VARIANTS} & 手順5で比べる構造（\file{["B", "A", "C"]}）と H の扱い（\file{["1", "2", "4", "14", "24"]}）。 & ― & 5 \\
\file{SPLITS} & 手順5で使う評価の分け方（版ごと。5data 版：\file{S_V}・\file{S_T}・\file{CV}、Alldata 版：\file{S_T}・\file{CV}・\file{S_Vall}）。 & ― & 5 \\
\file{E_GRID_A}、\file{EH_GRID}、\file{EV_GRID}、\file{EH_GRID_CV}、\file{EV_GRID_CV} & プロファイル尤度の粗い格子（A の $\Ev$、B の $\Eh$・$\Ev$。\file{_CV} は CV の各分割用の狭い格子）。 & kJ/mol & 5 \\
\file{FINE_A}、\file{FINE_B}、\file{FINE_B_CV} & 細かい格子の幅（最良の点の $\pm$。A は 0.5、B は 1 kJ/mol 刻み）。 & kJ/mol & 5 \\
\file{N_RESTARTS_PROFILE} & 尤度の格子の各点での最適化のやり直し回数（3）。 & ― & 5 \\
\file{B_EH_FIX} & None：B の $\Eh$・$\Ev$ を両方推定／数：$\Eh$ をその値に固定（$\Ev$ だけ推定）。 & kJ/mol & 5 \\
\file{P32_DRAWS} & (4) で \PII\ を予測して使うときの、予測分布からの引き直しの回数（20）。 & ― & 5 \\
\file{FIG_VARIANTS} & 曲線などの図に出す H の扱い（\file{["1", "2", "4"]}）。 & ― & 5 \\
\file{HR_VERSION}、\file{HR_MODELS} & 手順5の高解像度のグラフで描く版と、モデル（\file{"B_1"} など）。 & ― & 5 \\
\file{HR_...} & 高解像度のグラフの設定（レイアウト、帯、点、軸の範囲、文字の大きさ、保存形式など）。ノートの「高解像度のグラフ」のセルの先頭にある。 & ― & 3b・3c・3d・5a \\
\end{symtab}

\subsection{ノートの入力の列（ln は自然対数）}
\begin{symtab}
\file{ln_gd}、\file{ln_gd_cm} & $\ln\gd$（ショットごと）、$\ln\gd$ の条件平均。 & ― & 3a〜 \\
\file{ln_H} & $\ln H$。 & ― & 3a〜 \\
\file{ln_Tff}、\file{ln_Tff_cm} & $\ln\Tff$（K。ショットごと／条件平均）。 & ― & 3a〜 \\
\file{ln_P31}、\file{ln_P32}、\file{ln_P32_cm} & $\ln$\PI、$\ln$\PII（Pa に直す。ショットごと／条件平均）。 & ― & 3a〜 \\
\file{ln_v} & $\ln v$（射出速度）。 & ― & 3b〜 \\
\file{Tset} & $\Tset$（℃のまま）。 & ― & 3b〜 \\
\file{ln_eta} & $\ln\eta$（目的変数）。 & ― & 全 \\
\file{eta}・\file{gd}・\file{P31}・\file{P32}・\file{Tff}・\file{H}・\file{v}・\file{shot}・\file{cond}・\file{file} & CSV の列（それぞれ $\eta$・$\gd$・\PI・\PII（MPa）・$\Tff$（℃）・$H$・$v$・ショット番号・条件の名前・元ファイル名）。 & ― & 3a〜 \\
\end{symtab}

\subsection{結果の CSV の主な列（results の step3a・step3b・step3c\_mcmc・step5a）}
\begin{symtab}
\file{material}、\file{split}、\file{data}、\file{model}、\file{subset} & 材料、分け方、外れ値の扱い、モデル、評価する範囲（全体・低速側・高速側など）。 & ― & 3a〜 \\
\file{n_train}、\file{n_test}、\file{n_cond_test} & 学習・テストの行数、テストの条件数。 & ― & 3a〜 \\
\file{R2}、\file{RMSE}、\file{coverage95}、\file{NLPD} & 評価の指標。 & ― & 3a〜 \\
\file{pure_sd_ln_eta}、\file{pure_sd_after_gd} & 純粋誤差（同じ条件内の $\ln\eta$ の標準偏差）、そのうち $\ln\gd$ で説明した残り。 & ― & 3a・3b \\
\file{resid_sd}、\file{intercept} & 線形回帰の残差の標準偏差、切片。 & ― & 3a・3b \\
\file{ratio_190_230}、\file{observed} & 190/230 ℃ の粘度比（モデル／実測）。 & 倍 & 3b・5a \\
\file{structure}、\file{variant}、\file{set}、\file{mode} & 5a の構造（H・A・Hz・B ...）、同（metrics.csv）、尤度を求めた学習データ（S\_V・S\_T・ALL）、$\gd$ の入れ方（cm・shot）。 & ― & 5a \\
\file{E}、\file{Eh}、\file{Ev}、\file{LML}、\file{stage} & プロファイルの格子の $E$、$\Eh$、$\Ev$、対数周辺尤度、格子の段階（粗い・細かい）。 & ― & 5a \\
\file{E_best}、\file{Eh_best}、\file{Ev_best}、\file{LML_best} & 最良の $E$（B は $\Eh$・$\Ev$）と、そのときの LML。 & ― & 5a \\
\file{LML_at_E0}、\file{LML_at_E40}、\file{dLML_vs_E0} & $E=0$・$E=40$ での LML、$E=0$ との LML の差。 & ― & 5a \\
\file{E_95_lo}、\file{E_95_hi}、\file{at_grid_edge} & LML が最大から 1.92 以内の範囲、最良の $E$ が格子の端に張り付いたか。 & ― & 5a \\
\file{role} & main\_vs\_alt.csv の役割（本命・対案・基準・実測）。 & ― & 5a \\
\file{version} & 手順3d の版（\file{5data}・\file{Alldata}）。\file{versions.csv}・\file{metrics.csv}・\file{kernels.csv}・\file{summary.csv}。 & ― & 3d \\
\file{n_cond_train}、\file{shots_per_cond}、\file{gd_range}、\file{removed_by_list} & versions.csv：学習の条件数、1 条件のショット数の範囲、$\gd$ の範囲、問題データ一覧の「除外」で除いた行数。 & ― & 3d \\
\file{n_in_alldata}、\file{n_not_in_alldata} & inclusion.csv：その組の行のうち、Alldata に同じ行がある／ない行数（元ファイル名と、$\eta$・$\gd$・\PI・\PII・$\Tff$ が一致）。 & ― & 3d \\
\file{slope_per_shot}、\file{se}、\file{drift_first_to_last}、\file{z} & shot\_drift.csv：順番 \file{pos} に対する傾き（1 ショットあたり）、標準誤差、最初→最後のずれ（傾き × 順番の幅）、傾き / 標準誤差。 & ― & 3d \\
\file{half_diff}、\file{half_diff_se}、\file{half_diff_n_cond} & shot\_drift.csv（PP）：条件ごとの「\file{pos} 6〜10 の平均 − \file{pos} 1〜5 の平均」を、条件にわたって平均した値、その標準誤差、条件数。 & ln η など & 3d \\
\file{resid}、\file{sd}、\file{z}、\file{flag}、\file{list_id}、\file{data} & outlier\_check.csv：規則 f の残差（$\ln\eta$）、予測標準偏差、$z$、拾ったか、問題データ一覧の ID、データ（\file{全条件}：除外する前／\file{除外後}）。 & ― & 3d \\
\file{pure_sd_ln_eta}、\file{fit_sec} & kernels.csv：学習データの純粋誤差、学習にかかった秒数。 & ― & 3d・3e \\
\file{coef_*}、\file{RMSE_test}、\file{RMSE_over_pure} & formula\_check.csv（手順3e）：線形回帰の係数、検証の RMSE、その純粋誤差に対する比（リークならほぼ 0）。 & ― & 3e \\
\file{ratio_mean}、\file{ratio_min}、\file{ratio_max}、\file{gd_lo}、\file{gd_hi} & thickness\_ratio.csv：粘度の比（$\gd$ が重なる範囲での平均・最小・最大）と、その範囲。 & 倍、1/s & 3e \\
\file{bias} & bias\_by\_thickness.csv：厚みごとの残差の平均（実測 $-$ 予測、$\ln\eta$）。 & ln η & 3e \\
\file{R2_lnH_base}、\file{R2_lnH_with_P32}、\file{corr_lnH_lnP32}、\file{gain} & p32\_proxy.csv：$\ln H$ を条件平均の入力で回帰したときの R2（\PII\ なし／あり）、$\ln H$ と $\ln\PII$ の相関、R2 の増加。0.3 以上増えれば、\PII\ が厚みの代わりの目印になりうる。 & ― & 3e \\
\file{regression}、\file{coef_le}、\file{coef_lg}、\file{coef_lH}、\file{R2} & p32\_regression.csv：$\ln\PII$（比較は $\ln\PI$）を、条件平均の $\ln\eta$・$\ln\gd$・$\ln H$ で回帰した係数と R2（\file{le}・\file{lg}・\file{lH} は、それぞれ $\ln\eta$・$\ln\gd$・$\ln H$）。 & ― & 3e \\
\file{struct}、\file{hv}、\file{model}、\file{psrc} & 手順5 の構造（B・A・C）、H の扱い（\file{1}・\file{2}・\file{4}・\file{14}・\file{24}）、モデル名（\file{B_1} など）、\PII\ の出どころ（実測・予測・\file{―}＝\PII\ を使わない）。 & ― & 5 \\
\file{diff_best}、\file{diff_95_lo}、\file{diff_95_hi}、\file{n_in_95}、\file{at_edge} & best\_E.csv（手順5）：$\Ev-\Eh$ の最良の値と、LML が最大から 1.92 以内の範囲、その範囲の格子点の数、最良の点が格子の端に張り付いたか。 & kJ/mol & 5 \\
\file{amp_<m|c|f>}、\file{lin_amp_*}、\file{lin_sigma0_*}、\file{ls_<m|c|f>_<入力>} & kernels.csv（手順5）：部品ごとの振幅・線形項・長さスケール。\file{f}：(1)(4)(1+4) の 1 つのカーネル、\file{m}・\file{c}：(2)(2+4) の足し算の 2 つの成分。入力の記号は \file{u}・\file{H}・\file{T}・\file{w}・\file{p}。 & ― & 5 \\
\file{kind}（decomposition.csv）& \file{m_curve}：$m$ の曲線（設定温度ごと、$\ln\gd$ に対して）、\file{c_curve}：$c$ の曲線（$\ln\mathrm{Fo}$ に対して）、\file{c_cond}：学習の条件での $c$。 & ― & 5 \\
\file{actual}、\file{pred}、\file{pred_sd} & p32\_pred.csv（手順5）：検証の条件の $\ln\PII$（条件平均）の実測、予測、予測の標準偏差。 & ― & 5 \\
\file{chain}、\file{draw}、\file{lp}、\file{accepted} & MCMC のチェーンの番号、反復の番号、対数事後密度、採択したか（posterior\_*.csv）。 & ― & 3c \\
\end{symtab}

\subsection{高解像度のグラフのファイル名}
\file{results/<手順>/highres/hr_<材料>_<モデル>_<学習>_by-thickness_<厚み>mm.png}（厚みごと）と \file{..._by-temperature_<温度>C.png}（設定温度ごと）。3c は、\file{_mcmc} が付く。3d は、\file{hr_<材料>_R3_<版>_by-...}（版は \file{5data}・\file{Alldata}）。手順3d の図（\file{fig1}〜\file{fig5}）は、同じ名前で \file{highres/} にも 600 dpi で保存される（\file{HR_ALL}）。Git には入れない（ノートのセルで、いつでも作り直せる）。

\subsection{ノートの判定表示}
\file{>> 判定：OK} は、文の中に書いた基準を満たす（想定どおり）。\file{>> 判定：★要確認} は、満たさない（想定と違う、または結果を読むとき注意が要る）。エラーではない。基準は、判定の文の中に書いてある。

\end{document}
